\chapter{Advanced Prescriptive Analysis for multi-agent waterflooding optimization}
%applying actor-critic for multiple comparative-comparative agents on 

\section{Introduction}
Multi-Agent Actor-Critic for Mixed Cooperative-Competitive Environments
We explore deep reinforcement learning methods for multi-agent domains. We begin by analyzing the difficulty of traditional algorithms in the multi-agent case: Q-learning is challenged by an inherent non-stationarity of the environment, while policy gradient suffers from a variance that increases as the number of agents grows. We then present an adaptation of actor-critic methods that considers action policies of other agents and is able to successfully learn policies that require complex multi-agent coordination. Additionally, we introduce a training regimen utilizing an ensemble of policies for each agent that leads to more robust multi-agent policies. We show the strength of our approach compared to existing methods in cooperative as well as competitive scenarios, where agent populations are able to discover various physical and informational coordination strategies.

\section{DDPG}

In this post, I am sharing my understanding regarding Deterministic Policy Gradient Algorithm (DPG) [1] and its deep-learning version (DDPG) [2].

We have introduced policy gradient theorem in [3, 4]. Here, we briefly recap. The objective function of policy gradient methods is:

 J(\theta)=\sum\limits_{s \in S} d^\pi(s) V^\pi(s)=\sum\limits_{s \in S} d^\pi(s) \sum\limits_{a \in A} \pi(a|s) Q^\pi(s,a)

where $latex \pi$ represents $latex \pi_\theta$, $latex d^\pi(s)$ is the stationary distribution of Markov chain for $latex \pi$, $latex V^\pi(s)=\mathbb{E}_{a \sim \pi}[G_t|S_t=s]$, and $latex Q^{\pi}(s,a)=\mathbb{E}_{a \sim \pi}[G_t|S_t=s, A_t=a]$. $latex G_t$ is accumulated rewards since time step $latex t$: $latex G_t=\sum^\infty_{k=0}\gamma^k R_{t+k}$.

Policy gradient theorem proves that the gradient of policy parameters with regard to $ J(\theta)$ is:

\begin{equation}
    \nabla_\theta J(\theta) = \mathbb{E}_{\pi}[Q^\pi(s,a)\nabla_\theta ln \pi_\theta(a|s)] &s=2
\end{equation}

More specifically, the policy gradient theorem we are talking about is stochastic policy gradient theorem because at each state the policy outputs a stochastic action distribution $\pi(s|a)$.

However, the policy gradient theorem implies that policy gradient must be calculated on-policy, as $\mathbb{E}_{\pi}:=\mathbb{E}_{s\sim d^{\pi}, a\sim \pi}$ means the state and action distribution is generated by following $latex \pi_\theta$.

If the training samples is generated by some other behavioral policy $latex \beta$, then the objective function becomes:

\begin{equation}
    J(\theta)=\sum\limits_{s \in S} d^\beta(s) V^\pi(s)=\sum\limits_{s \in S} d^\beta(s) \sum\limits_{a \in A} \pi(a|s) Q^\pi(s,a), &s=2
\end{equation}

and we must rely on importance sampling to calculate the policy gradient [7,8]:
\begin{equation}
    \nabla_\theta J(\theta) = \mathbb{E}_{s\sim d^\beta}[\sum\limits_{a\in A}\nabla_\theta \pi(a|s) Q^\pi(s,a) + \pi(a|s)\nabla_\theta Q^\pi(s,a)] \newline \approx \mathbb{E}_{s\sim d^\beta}[\sum\limits_{a\in A}\nabla_\theta \pi(a|s) Q^\pi(s,a)] \newline =\mathbb{E}_{\beta}[\frac{\pi_\theta(a|s)}{\beta(a|s)}Q^\pi(s,a)\nabla_\theta ln \pi_\theta(a|s)] &s=2    
\end{equation}

where $\mathbb{E}_{\beta}:=\mathbb{E}_{s\sim d^\beta, a\sim\beta}$

However, using importance sampling $ \frac{\pi_\theta(a|s)}{\beta(a|s)}$ could incur some learning difficulty problems:


What DPG proposes is that we design a policy that outputs actions deterministically: $latex a = \mu_\theta(s)$, thus $latex J(\theta)$ would only have integration (or sum) over state distribution and get rid of importance sampling, which makes learning potentially easier:

$latex J(\theta)=\sum\limits_{s \in S} d^\beta(s) V^\mu(s)=\sum\limits_{s \in S} d^\beta(s) Q^\mu(s, \mu_\theta(s)) &s=2$

Writing the above formula in integration (rather than discrete sum), we get:

$latex J(\theta)=\int\limits_S d^\beta(s) V^\mu(s) ds=\int\limits_{S} d^\beta(s) Q^\mu(s, \mu_\theta(s)) ds &s=2$

And the gradient of $latex J(\theta)$ is:

$latex \nabla_\theta J(\theta) = \int\limits_S d^\beta(s) \nabla_\theta Q^\mu(s, \mu_\theta(s)) ds \approx \mathbb{E}_\beta [\nabla_\theta \mu_\theta(s) \nabla_a Q^\mu(s,a)|_{a=\mu_\theta(s)}] &s=2$

This immediately implies that $latex \theta$ can be updated off policy. Because we only need to sample $latex s \sim d^\beta(s)$, then $latex \nabla J(\theta)$ can be calculated without knowing what action the behavior policy $latex \beta$ took (we only need to know what the training policy would take, i.e., $latex a=\mu_\theta(s)$). We should also note that there is an approximation in the formula of $latex \nabla_\theta J(\theta)$, because $latex \nabla_\theta Q^u(s,\mu_\theta(s))&s=2$ cannot be applied chain rules easily since even you replace $latex \mu_\theta(s))$ with $latex a$, the $latex Q$ function itself still depends on $latex \theta$.

Updating $latex \theta$ is one part of DPG. The other part is fitting a Q-network with parameter $latex \phi$ on the optimal Q-function of policy governed by $latex \theta$. This equates to update according to mean-squared Bellman error (MSBE) [10], which is the same as in Q-Learning.

I haven’t read the DDPG paper [2] thoroughly but based on my rough understanding, it adds several tricks when dealing with large inputs (such as images). According to [9]’s summary, DDPG introduced three tricks:

add batch normalization to normalize “every dimension across samples in one minibatch”
add noise in the output of DPG such that the algorithm can explore
soft update target network


\section{2}

Deep Deterministic Policy Gradient (DDPG) is an algorithm which
concurrently learns a Q-function and a policy. It uses off-policy data
and the Bellman equation to learn the Q-function, and uses the
Q-function to learn the policy.

This approach is closely connected to Q-learning, and is motivated the
same way: if you know the optimal action-value function \(Q^*(s,a)\),
then in any given state, the optimal action \(a^*(s)\) can be found by
solving

\[a^*(s) = \arg \max_a Q^*(s,a).\]

DDPG interleaves learning an approximator to \(Q^*(s,a)\) with learning
an approximator to \(a^*(s)\), and it does so in a way which is
specifically adapted for environments with continuous action spaces. But
what does it mean that DDPG is adapted \emph{specifically} for
environments with continuous action spaces? It relates to how we compute
the max over actions in \(\max_a Q^*(s,a)\).

When there are a finite number of discrete actions, the max poses no
problem, because we can just compute the Q-values for each action
separately and directly compare them. (This also immediately gives us
the action which maximizes the Q-value.) But when the action space is
continuous, we can't exhaustively evaluate the space, and solving the
optimization problem is highly non-trivial. Using a normal optimization
algorithm would make calculating \(\max_a Q^*(s,a)\) a painfully
expensive subroutine. And since it would need to be run every time the
agent wants to take an action in the environment, this is unacceptable.

Because the action space is continuous, the function \(Q^*(s,a)\) is
presumed to be differentiable with respect to the action argument. This
allows us to set up an efficient, gradient-based learning rule for a
policy \(\mu(s)\) which exploits that fact. Then, instead of running an
expensive optimization subroutine each time we wish to compute
\(\max_a Q(s,a)\), we can approximate it with
\(\max_a Q(s,a) \approx Q(s,\mu(s))\). See the Key Equations section
details.

----------------------

First, let's recap the Bellman equation describing the optimal
action-value function, \(Q^*(s,a)\). It's given by

\[Q^*(s,a) = \underset{s' \sim P}{{\mathrm E}}\left[r(s,a) + \gamma \max_{a'} Q^*(s', a')\right]\]

where \(s' \sim P\) is shorthand for saying that the next state, \(s'\),
is sampled by the environment from a distribution \(P(\cdot| s,a)\).

This Bellman equation is the starting point for learning an approximator
to \(Q^*(s,a)\). Suppose the approximator is a neural network
\(Q_{\phi}(s,a)\), with parameters \(\phi\), and that we have collected
a set \({\mathcal D}\) of transitions \((s,a,r,s',d)\) (where \(d\)
indicates whether state \(s'\) is terminal). We can set up a
\textbf{mean-squared Bellman error (MSBE)} function, which tells us
roughly how closely \(Q_{\phi}\) comes to satisfying the Bellman
equation:

\[L(\phi, {\mathcal D}) = \underset{(s,a,r,s',d) \sim {\mathcal D}}{{\mathrm E}}\left[
    \Bigg( Q_{\phi}(s,a) - \left(r + \gamma (1 - d) \max_{a'} Q_{\phi}(s',a') \right) \Bigg)^2
    \right]\]

Here, in evaluating \((1-d)\), we've used a Python convention of
evaluating \texttt{True} to 1 and \texttt{False} to zero. Thus, when
\texttt{d==True}-\/-\/-which is to say, when \(s'\) is a terminal
state-\/-\/-the Q-function should show that the agent gets no additional
rewards after the current state. (This choice of notation corresponds to
what we later implement in code.)

Q-learning algorithms for function approximators, such as DQN (and all
its variants) and DDPG, are largely based on minimizing this MSBE loss
function. There are two main tricks employed by all of them which are
worth describing, and then a specific detail for DDPG.

\textbf{Trick One: Replay Buffers.} All standard algorithms for training
a deep neural network to approximate \(Q^*(s,a)\) make use of an
experience replay buffer. This is the set \({\mathcal D}\) of previous
experiences. In order for the algorithm to have stable behavior, the
replay buffer should be large enough to contain a wide range of
experiences, but it may not always be good to keep everything. If you
only use the very-most recent data, you will overfit to that and things
will break; if you use too much experience, you may slow down your
learning. This may take some tuning to get right.

You Should Know

We've mentioned that DDPG is an off-policy algorithm: this is as good a
point as any to highlight why and how. Observe that the replay buffer
\emph{should} contain old experiences, even though they might have been
obtained using an outdated policy. Why are we able to use these at all?
The reason is that the Bellman equation \emph{doesn't care} which
transition tuples are used, or how the actions were selected, or what
happens after a given transition, because the optimal Q-function should
satisfy the Bellman equation for \emph{all} possible transitions. So any
transitions that we've ever experienced are fair game when trying to fit
a Q-function approximator via MSBE minimization.

\textbf{Trick Two: Target Networks.} Q-learning algorithms make use of
\textbf{target networks}. The term

\[r + \gamma (1 - d) \max_{a'} Q_{\phi}(s',a')\]

is called the \textbf{target}, because when we minimize the MSBE loss,
we are trying to make the Q-function be more like this target.
Problematically, the target depends on the same parameters we are trying
to train: \(\phi\). This makes MSBE minimization unstable. The solution
is to use a set of parameters which comes close to \(\phi\), but with a
time delay-\/-\/-that is to say, a second network, called the target
network, which lags the first. The parameters of the target network are
denoted \(\phi_{\text{targ}}\).

In DQN-based algorithms, the target network is just copied over from the
main network every some-fixed-number of steps. In DDPG-style algorithms,
the target network is updated once per main network update by polyak
averaging:

\[\phi_{\text{targ}} \leftarrow \rho \phi_{\text{targ}} + (1 - \rho) \phi,\]

where \(\rho\) is a hyperparameter between 0 and 1 (usually close to 1).
(This hyperparameter is called \texttt{polyak} in our code).

\textbf{DDPG Detail: Calculating the Max Over Actions in the Target.} As
mentioned earlier: computing the maximum over actions in the target is a
challenge in continuous action spaces. DDPG deals with this by using a
\textbf{target policy network} to compute an action which approximately
maximizes \(Q_{\phi_{\text{targ}}}\). The target policy network is found
the same way as the target Q-function: by polyak averaging the policy
parameters over the course of training.

Putting it all together, Q-learning in DDPG is performed by minimizing
the following MSBE loss with stochastic gradient descent:

\[L(\phi, {\mathcal D}) = \underset{(s,a,r,s',d) \sim {\mathcal D}}{{\mathrm E}}\left[
    \Bigg( Q_{\phi}(s,a) - \left(r + \gamma (1 - d) Q_{\phi_{\text{targ}}}(s', \mu_{\theta_{\text{targ}}}(s')) \right) \Bigg)^2
    \right],\]

where \(\mu_{\theta_{\text{targ}}}\) is the target policy.

\hypertarget{the-policy-learning-side-of-ddpg}{%
\subsubsection{The Policy Learning Side of
DDPG}\label{the-policy-learning-side-of-ddpg}}

Policy learning in DDPG is fairly simple. We want to learn a
deterministic policy \(\mu_{\theta}(s)\) which gives the action that
maximizes \(Q_{\phi}(s,a)\). Because the action space is continuous, and
we assume the Q-function is differentiable with respect to action, we
can just perform gradient ascent (with respect to policy parameters
only) to solve

\[\max_{\theta} \underset{s \sim {\mathcal D}}{{\mathrm E}}\left[ Q_{\phi}(s, \mu_{\theta}(s)) \right].\]

Note that the Q-function parameters are treated as constants here.

\hypertarget{exploration-vs.-exploitation}{%
\subsection{Exploration vs.
Exploitation}\label{exploration-vs.-exploitation}}

DDPG trains a deterministic policy in an off-policy way. Because the
policy is deterministic, if the agent were to explore on-policy, in the
beginning it would probably not try a wide enough variety of actions to
find useful learning signals. To make DDPG policies explore better, we
add noise to their actions at training time. The authors of the original
DDPG paper recommended time-correlated
\href{https://en.wikipedia.org/wiki/Ornstein\%E2\%80\%93Uhlenbeck_process}{OU
noise}, but more recent results suggest that uncorrelated, mean-zero
Gaussian noise works perfectly well. Since the latter is simpler, it is
preferred. To facilitate getting higher-quality training data, you may
reduce the scale of the noise over the course of training. (We do not do
this in our implementation, and keep noise scale fixed throughout.)

At test time, to see how well the policy exploits what it has learned,
we do not add noise to the actions.

You Should Know

Our DDPG implementation uses a trick to improve exploration at the start
of training. For a fixed number of steps at the beginning (set with the
\texttt{start\_steps} keyword argument), the agent takes actions which
are sampled from a uniform random distribution over valid actions. After
that, it returns to normal DDPG exploration.

\begin{algorithm}[H]
    \caption{Deep Deterministic Policy Gradient}
    \label{alg1}
\begin{algorithmic}[1]
    \STATE Input: initial policy parameters $\theta$, Q-function parameters $\phi$, empty replay buffer $\mathcal{D}$
    \STATE Set target parameters equal to main parameters $\theta_{\text{targ}} \leftarrow \theta$, $\phi_{\text{targ}} \leftarrow \phi$
    \REPEAT
        \STATE Observe state $s$ and select action $a = \text{clip}(\mu_{\theta}(s) + \epsilon, a_{Low}, a_{High})$, where $\epsilon \sim \mathcal{N}$
        \STATE Execute $a$ in the environment
        \STATE Observe next state $s'$, reward $r$, and done signal $d$ to indicate whether $s'$ is terminal
        \STATE Store $(s,a,r,s',d)$ in replay buffer $\mathcal{D}$
        \STATE If $s'$ is terminal, reset environment state.
        \IF{it's time to update}
            \FOR{however many updates}
                \STATE Randomly sample a batch of transitions, $B = \{ (s,a,r,s',d) \}$ from $\mathcal{D}$
                \STATE Compute targets
                \begin{equation*}
                    y(r,s',d) = r + \gamma (1-d) Q_{\phi_{\text{targ}}}(s', \mu_{\theta_{\text{targ}}}(s'))
                \end{equation*}
                \STATE Update Q-function by one step of gradient descent using
                \begin{equation*}
                    \nabla_{\phi} \frac{1}{|B|}\sum_{(s,a,r,s',d) \in B} \left( Q_{\phi}(s,a) - y(r,s',d) \right)^2
                \end{equation*}
                \STATE Update policy by one step of gradient ascent using
                \begin{equation*}
                    \nabla_{\theta} \frac{1}{|B|}\sum_{s \in B}Q_{\phi}(s, \mu_{\theta}(s))
                \end{equation*}
                \STATE Update target networks with
                \begin{align*}
                    \phi_{\text{targ}} &\leftarrow \rho \phi_{\text{targ}} + (1-\rho) \phi \\
                    \theta_{\text{targ}} &\leftarrow \rho \theta_{\text{targ}} + (1-\rho) \theta
                \end{align*}
            \ENDFOR
        \ENDIF
    \UNTIL{convergence}
\end{algorithmic}
\end{algorithm}


\section{MDDPG}
Multi-Agent Actor Critic

Figure 1: Overview of our multi-agent decentralized actor, centralized critic approach.
We have argued in the previous section that naïve
policy gradient methods perform poorly in simple
multi-agent settings, and this is supported in our experiments in Section 5. Our goal in this section is to
derive an algorithm that works well in such settings.
However, we would like to operate under the following constraints: (1) the learned policies can only use
local information (i.e. their own observations) at execution time, (2) we do not assume a differentiable
model of the environment dynamics, unlike in [25],
and (3) we do not assume any particular structure on
the communication method between agents (that is, we don’t assume a differentiable communication
channel). Fulfilling the above desiderata would provide a general-purpose multi-agent learning
algorithm that could be applied not just to cooperative games with explicit communication channels,
but competitive games and games involving only physical interactions between agents.
Similarly to [8], we accomplish our goal by adopting the framework of centralized training with
decentralized execution. Thus, we allow the policies to use extra information to ease training, so
long as this information is not used at test time. It is unnatural to do this with Q-learning, as the Q
function generally cannot contain different information at training and test time. Thus, we propose
a simple extension of actor-critic policy gradient methods where the critic is augmented with extra
information about the policies of other agents.
More concretely, consider a game with N agents with policies parameterized by θ = {θ1, ..., θN },
and let π = {π1, ..., πN } be the set of all agent policies. Then we can write the gradient of the
expected return for agent i, J(θi) = E[Ri
] as:

----------------
\section{Methods}
\label{sec:methods}
\subsection{Multi-Agent Actor Critic}\label{sec:maac}

\begin{figure}{r}{0.45\textwidth}
\includegraphics[width=.9\linewidth]{figures/overview_compact}
\caption{\label{fig:model} Overview of our multi-agent decentralized actor, centralized critic approach}
\end{figure}

We have argued in the previous section %%PA: little strange to start off our own work with a claim of justification that's in the appendix, maybe just say (and our experiments also support this)
that na{\"i}ve policy gradient methods perform poorly in simple multi-agent settings, and this is supported in our experiments in Section \ref{sec:experiments}. Our goal in this section is to derive an algorithm that works well %%PA: empirically --> well ?
in such settings. However, we would like to operate under the following constraints: (1) the learned policies can only use local information (i.e.\@ their own observations) at execution time, (2) we do not assume a differentiable model of the environment dynamics, unlike in \cite{mordatch2017emergence}, %%PA: it's unclear what "as in" means -- does [23] use diff assumptions, or do we just like [23] not assume it (maybe just cut this reference here?)
and (3) we do not assume any particular structure on the communication method between agents (that is, we don't assume a differentiable communication channel). Fulfilling the above desiderata would provide a general-purpose multi-agent learning algorithm that could be applied not just to cooperative games with explicit communication channels, but competitive games and games involving only physical interactions between agents.

%\insertfigure{overview_compact}{0.40}{Overview of our multi-agent decentralized actor, centralized critic approach.}

Similarly to \cite{foerster16b}, we accomplish our goal by adopting the framework of centralized training with decentralized execution. Thus, we allow the policies to use extra information to ease training, so long as this information is not used at test time. It is unnatural to do this with Q-learning, as the Q function generally cannot contain different information at training and test time. Thus, we propose a simple extension of actor-critic policy gradient methods where the critic is augmented with extra information about the policies of other agents. %Agents then increase the probability of executing a certain action


More concretely, consider a game with $N$ agents with policies parameterized by $\pmb{\theta} = \{\theta_1, ..., \theta_N\}$, and let $\pol = \{\pol_1, ..., \pol_N\}$ be the set of all agent policies. Then we can write the gradient of the expected return for agent $i$, $J(\theta_i) = \mathbb{E}[R_i]$ as:
\begin{equation}
\nabla_{\theta_i} J(\theta_i) = \mathbb{E}_{s\sim p^{\cpol},a_i \sim \pol_i} [\nabla_{\theta_i} \log \pol_i(a_i|o_i) Q^{\pol}_i (\mathbf{x},a_1, ..., a_N)]. 
\end{equation}

Here $Q^{\pol}_i (\mathbf{x},a_1, ..., a_N)$ is a \textit{centralized action-value function} that takes as input the actions of all agents, $a_1,\ldots, a_N$, in addition to some state information $\mathbf{x}$, and outputs the Q-value for agent $i$. In the simplest case, $\mathbf{x}$ could consist of the observations of all agents, $\mathbf{x} = (o_1, ..., o_N)$, however we could also include additional state information if available. Since each $Q^{\pol}_i$ is learned separately, agents can have arbitrary reward structures, including conflicting rewards in a competitive setting.

We can extend the above idea to work with deterministic policies. If we now consider $N$ continuous policies $\cpol_{\theta_i}$ w.r.t. parameters $\theta_i$ (abbreviated as $\cpol_i$), the gradient can be written as:
\begin{equation}
\nabla_{\theta_i} J(\cpol_i) = \mathbb{E}_{\mathbf{x},a \sim \mathcal{D}}[\nabla_{\theta_i} \cpol_i(a_i|o_i) \nabla_{a_i} Q^{\cpol}_i (\mathbf{x},a_1, ..., a_N)|_{a_i=\cpol_i (o_i)}],
\end{equation}
Here the experience replay buffer $\mathcal{D}$ contains the tuples $(\mathbf{x},\mathbf{x}',a_1,\ldots,a_N,r_1,\ldots,r_N)$, recording experiences of all agents. 
The centralized action-value function $Q^{\cpol}_i$ is updated as:
\begin{equation}\label{eq:q_func}
\mathcal{L}(\theta_i) = \mathbb{E}_{\mathbf{x},a,r,\mathbf{x}'}[(Q^{\cpol}_i(\mathbf{x},a_1,\ldots,a_N) - y)^2], \;\;\;\;
y = r_i + \gamma\, Q^{\cpol'}_i(\mathbf{x}', a_1',\ldots,a_N')\big|_{a_j'=\cpol'_j(o_j)},
%\label{eq:maddpg_target_y} Yi: to save some space, I merged these two lines
\end{equation}
where $\cpol' = \{\cpol_{\theta'_1}, ..., \cpol_{\theta'_N} \}$ is the set of target policies with delayed parameters $\theta'_i$. As shown in Section \ref{sec:experiments}, we find the centralized critic with deterministic policies works very well in practice, and refer to it as \textit{multi-agent deep deterministic policy gradient} (MADDPG). 
We provide the description of the full algorithm in the Appendix. 

A primary motivation behind MADDPG is that, if we know the actions taken by all agents, the environment is stationary even as the policies change, since $P(s'|s,a_1,...,a_N,\pol_1,...,\pol_N) = P(s'|s,a_1,...,a_N) = P(s'|s,a_1,...,a_N,\pol'_1,...,\pol'_N)$ for any $\pol_i \neq \pol_i'$. This is not the case if we do not explicitly condition on the actions of other agents, as done for most traditional RL methods. %, as for traditional RL methods, as seen in Section \ref{sec:background}.
% TODO: add below
%In addition, by conditioning on the actions of other agents we are reducing the variance of the gradients to each agent; the centralized critic is able to 

Note that we require the policies of other agents to apply an update in Eq.~\ref{eq:q_func}. Knowing the observations and policies of other agents is not a particularly restrictive assumption; if our goal is to train agents to exhibit complex communicative behaviour in simulation, this information is often available to all agents.
However, we can relax this assumption if necessary by learning the policies of other agents from observations --- we describe a method of doing this in Section \ref{sec:modellearning}.


\subsection{Inferring Policies of Other Agents}
\label{sec:modellearning}
%\comment{Yi: This section should be modified after sec4.1 is updated.}
%In our algorithm framework, when updating the Q function of an agent (eq.~\ref{eq:q_func}), we need to compute the action of other agents using their true policies. This means, each agent's policy function must be fully  accessible for other agents during training. 
%Actually, we can even get rid of this requirement by approximating other agents policy functions.
To remove the assumption of knowing other agents' policies, as required in Eq.~\ref{eq:q_func}, each agent $i$ can additionally maintain an approximation $\hat{\cpol}_{\phi_i^j}$ (where $\phi$ are the parameters of the approximation; henceforth $\hat{\cpol}_i^j$) to the true policy of agent $j$, $\cpol_j$. This approximate policy is learned by maximizing the log probability of agent $j$'s actions, with an entropy regularizer:
%For agent $i$, besides its own policy function $\cpol_{\theta_i}$, we also maintains an approximate policy $\hat{\cpol}_{\phi_i^j}$ to approximate the true policy of agent $j$, $\cpol_{\theta_j}$ by minimizing the loss function
\begin{equation}\label{eq:loss-modellearn}
\mathcal{L}(\phi_i^j)=-\mathbb{E}_{o_j,a_j}\left[\log \hat{\cpol}_i^j(a_j|o_j) + \lambda H(\hat{\cpol}_i^j) \right],
\end{equation}
where $H$ is the entropy of the policy distribution. With the approximate policies, $y$ in Eq.~\ref{eq:q_func} can be replaced by an approximate value $\hat{y}$ calculated as follows: 
\begin{equation}\label{eq:approx_q_func}
\hat{y} = r_i + \gamma Q^{\cpol'}_i(\mathbf{x}', \hat{\cpol}'^1_i(o_1),\ldots,\cpol'_i(o_i),\ldots,\hat{\cpol}'^N_i(o_N)),
\end{equation}
where $\hat{\cpol}'^j_i$ denotes the target network for the approximate policy $\hat{\cpol}_i^j$.
Note that Eq.~\ref{eq:loss-modellearn} can be optimized in a completely online fashion: before updating  $Q^{\cpol}_i$, the centralized Q function, we take the latest samples of each agent $j$ from the replay buffer to perform a single gradient step to update $\phi_i^j$. %While this formulation requires agent $j$'s observations, it is most pertinent in the setting where agents $i$ and $j$ have similar 
Note also that, in the above equation, we input the action log probabilities of each agent directly into $Q$, rather than sampling.

\subsection{Agents with Policy Ensembles}
\label{sec:ensemble}
As previously mentioned, a recurring problem in multi-agent reinforcement learning is the environment non-stationarity due to the agents' changing policies. This is particularly true in competitive settings, where agents can derive a strong policy by overfitting to the behavior of their competitors. Such policies are undesirable as they are brittle and may fail when the competitors alter strategies.

%In multi-agent reinforcement learning, the environment is non-stationary due to the agents' changing policies. When one agent's policy converges to a bad local optimum, this bad mode can accordingly affect other agents' policies. Especially in the adversary/competing scenarios, one agent can simply derive a strong policy by overfitting the competitor's behavior, which can simply fail when facing another slightly different strategy. \comment{Yi: I feel that I need a better example.... I just put one here. Feeling not very appropriate...}For example, in the rock-paper-scissor game, one stationary strategy is that both agents always raise rock (or paper/scissor), which is brittle since it will be always defeated by raising paper. The optimal strategy should be a random policy, which can be seen as a random combination of the 3 aforementioned stationary policies. 

To obtain multi-agent policies that are more robust to changes in the policy of competing agents, we propose to train a collection of $K$ different sub-policies. At each episode, we randomly select one particular sub-policy for each agent to execute. 
Suppose that policy $\cpol_i$ is an ensemble of $K$ different sub-policies with sub-policy $k$ denoted by $\cpol_{\theta_i^{(k)}}$ (denoted as $\cpol_i^{(k)}$). For agent $i$, we are then maximizing the ensemble objective:
%\begin{equation}\label{eq:J_ensemble}
$
J_e(\cpol_i)=\mathbb{E}_{k\sim\textrm{unif}(1,K),s\sim p^{\cpol},a\sim\cpol_i^{(k)}}\left[R_i(s,a)\right].
$
%\end{equation}

Since different sub-policies will be executed in different episodes, we maintain a replay buffer $\mathcal{D}_i^{(k)}$ for each sub-policy $\cpol_i^{(k)}$ of agent $i$. Accordingly, we can derive the gradient of the ensemble objective with respect to $\theta_i^{(k)}$ as follows:
\begin{equation}\label{eq:J_ensemble_grad}
\nabla_{\theta_i^{(k)}} J_e(\cpol_i)=\frac{1}{K}\mathbb{E}_{\mathbf{x},a\sim \mathcal{D}_i^{(k)}}\left[\nabla_{\theta_i^{(k)}} \cpol_i^{(k)}(a_i|o_i)\nabla_{a_i}Q^{\cpol_i}\left(\mathbf{x},a_1,\ldots,a_N\right)\Big|_{a_i=\cpol_i^{(k)}(o_i)}\right].
\end{equation}
